\Needspace{12\baselineskip}
\begin{center}
\textit{Forms with positive nonsquare determinant.}
\end{center}

\Needspace{8\baselineskip}
\Article{183}
\textsc{Problem.} \textit{Given any form $(a,b,a')$ with positive nonsquare determinant $=D$, find a properly equivalent form $(A,B,C)$ in which $B$ is positive and $<\surd D$, and $A$, if positive, or $-A$, if $A$ is negative, lies between $\surd D+B$ and $\surd D-B$.}

\textit{Solution.} We assume the given form does not yet satisfy both conditions; otherwise there would be no need to look for another form. We also note that in a form with \textit{nonsquare} determinant, the first or last term cannot be $=0$ (Article 171, note). Let $b'\equiv-b\ (\mathrm{mod.}\ a')$, lying between the bounds $\surd D$ and $\surd D\mp a'$ (taking the upper sign when $a'$ is positive and the lower when it is negative). That this is possible is easily proved as in Article 3. Put $\frac{b'b'-D}{a'}=a''$, which is an integer because $b'b'-D\equiv bb-D\equiv aa'\equiv0\ (\mathrm{mod.}\ a')$. Now if $a''<a'$, again let $b''\equiv-b'\ (\mathrm{mod.}\ a'')$, lying between $\surd D$ and $\surd D\mp a''$ (according as $a''$ is positive or negative), and $\frac{b''b''-D}{{a''}}=a'''$. If again $a'''<a''$, once more let $b'''\equiv-b''\ (\mathrm{mod.}\ a''')$, lying between $\surd D$ and $\surd D\mp a'''$, and $\frac{b'''b'''-D}{a'''}=a''''$. Continue this operation until the sequence $a',a'',a''',a''''$ etc. reaches a term $a^{m+1}$ no smaller than its predecessor $a^m$. This must eventually happen, since otherwise we would have an infinite sequence of integers that keep decreasing. Then, on putting $a^m=A$, $b^m=B$, $a^{m+1}=C$, the form $(A,B,C)$ satisfies all the conditions.

\textit{Proof.} I. In the sequence of forms $(a,b,a')$, $(a',b',a'')$, $(a'',b'',a''')$ etc., each is a neighbor of its predecessor; so the last, $(A,B,C)$, is properly equivalent to the first, $(a,b,a')$.

II. Since $B$ lies between $\surd D$ and $\surd D\mp A$ (always taking the upper sign when $A$ is positive and the lower when $A$ is negative), putting $\surd D-B=p$, $B-(\surd D\mp A)=q$ clearly makes $p,q$ positive. It is easy to confirm that $qq+2pq+2p\surd D=D+AA-BB$; so $D+AA-BB$ is a positive number, which we put $=r$. Then, since $D=BB-AC$,
\SourceLeaf{169}
we get $r=AA-AC$, so $AA-AC$ is positive. But by assumption $A$ is no greater than $C$, so this clearly requires $AC$ to be negative, making the signs of $A,C$ opposite. This gives $BB=D+AC<D$, and so $B<\surd D$.

III. Further, since $-AC=D-BB$, we have $AC<D$, and so (since $A$ is not $>C$) $A<\surd D$. Thus $\surd D\mp A$ is positive, as is $B$, which lies between $\surd D$ and $\surd D\mp A$.

IV. So $\surd D+B\mp A$ is certainly positive. Since $\surd D-B\mp A=-q$ is negative, $\pm A$ lies between $\surd D+B$ and $\surd D-B$. \textit{Q. E. D.}

\textit{Example.} Let the given form be $(67,97,140)$, whose determinant $=29$. Here we find the sequence of forms $(67,97,140)$, $(140,-97,67)$, $(67,-37,20)$, $(20,-3,-1)$, $(-1,5,4)$. The last is the one we want.

We will call forms $(A,B,C)$ with positive nonsquare determinant $D$, in which $A$, taken positively, lies between $\surd D+B$ and $\surd D-B$, while $B$ is positive and $<\surd D$, \textit{reduced forms}. Reduced forms with positive nonsquare determinant thus differ somewhat from reduced forms with negative determinant. But because of the close analogy between them, we did not want to introduce different names.

\Needspace{8\baselineskip}
\Article{184}
If we could recognize the equivalence of two \textit{reduced} forms with positive determinant as easily as for forms with negative determinant (Article 172), we could decide without difficulty whether two \textit{arbitrary} forms with the same positive determinant are equivalent. But the situation here is quite different: many reduced forms can be equivalent to one another. Before taking up this problem, we must look more deeply into the nature of reduced forms (with positive nonsquare determinant, which is always understood here).

1) If $(a,b,c)$ is a reduced form, $a$ and $c$ have opposite signs. For putting the determinant of the form $=D$, we have $ac=bb-D$, which is negative because $b<\surd D$.

2) The number $c$, like $a$, taken positively, lies between $\surd D+b$ and $\surd D-b$. For $-c=\frac{D-bb}{a}$, so, disregarding the sign, $c$ lies between $\frac{D-bb}{\surd D+b}$ and $\frac{D-bb}{\surd D-b}$, \textit{that is}, between $\surd D-b$ and $\surd D+b$.

3) So $(c,b,a)$ is also a reduced form.

\SourceLeaf{170}
4) Both $a$ and $c$ are $<2\surd D$. Each is $<\surd D+b$, so also $<2\surd D$.

5) The number $b$ lies between $\surd D$ and $\surd D\mp a$ (taking the upper sign when $a$ is positive and the lower when it is negative). Since $\pm a$ lies between $\surd D+b$ and $\surd D-b$, $\pm a-(\surd D-b)$, or $b-(\surd D\mp a)$, is positive. But $b-\surd D$ is negative, so $b$ lies between $\surd D$ and $\surd D\mp a$. --- Exactly the same argument shows that $b$ lies between $\surd D$ and $\surd D\mp c$ (according as $c$ is positive or negative).

6) \textit{Every reduced form $(a,b,c)$ has exactly one reduced neighbor on each side.}

Put $a'=c$, $b'\equiv-b\ (\mathrm{mod.}\ a')$, lying between $\surd D$ and $\surd D\mp a'$\HistoricalNote{*)}{Where there are alternative signs, the upper sign always applies when $a'$ is positive and the lower when $a'$ is negative.}, and $c'=\frac{b'b'-D}{a'}$. Then the form $(a',b',c')$ is a right neighbor of the form $(a,b,c)$. It is also clear that any reduced right neighbor of the form $(a,b,c)$ must be this $(a',b',c')$. That it really is reduced is proved as follows.

A) On putting
\[
\surd D+b\mp a'=p,\quad\pm a'-(\surd D-b)=q,\quad\surd D-b=r
\]
these $p,q,r$ are positive, by (2) above and the definition of a reduced form. Further, put
\[
b'-(\surd D\mp a)=q',\quad\surd D-b'=r'
\]
Then $q',r'$ are positive, since $b'$ lies between $\surd D$ and $\surd D\mp a'$. Finally, let $b+b'=\pm ma$; then $m$ is an integer. Clearly $p+q'=b+b'$, so $b+b'$, or $\pm ma'$, is positive, and therefore so is $m$. Thus $m-1$ is certainly nonnegative. Further, we get
\[
r+q'\pm ma'=2b'\pm a',\quad\text{or}\quad2b'=r+q'\pm(m-1)a'
\]
so $2b'$ and $b'$ must be positive. Since $b'+r'=\surd D$, we have $b'<\surd D$.

\Needspace{9\baselineskip}
B) Further, we get
\[
r\pm ma'=\surd D+b',\quad\text{or}\quad r\pm(m-1)a'=\surd D+b'\mp a'
\]
so $\surd D+b'\mp a'$ is positive. From this, and since $\pm a'-(\surd D-b')=q'$,
\SourceLeaf{171}
and so is positive, $\pm a'$ lies between $\surd D+b'$ and $\surd D-b'$. --- Thus $(a',b',c')$ is a reduced form.

The same argument proves that if we put ${}'c=a$, ${}'b\equiv-b\ (\mathrm{mod.}\ {}'c)$, lying between $\surd D$ and $\surd D\pm{}'c$, and ${}'a=\frac{{}'b{}'b-D}{{}'c}$, the form $({}'a,{}'b,{}'c)$ is reduced. Clearly this form is a left neighbor of the form $(a,b,c)$, and no reduced form other than $({}'a,{}'b,{}'c)$ can have this property.

\textit{Example.} The reduced form $(5,11,-14)$, whose determinant $=191$, has the reduced right neighbor $(-14,3,13)$ and the left neighbor $(-22,9,5)$.

7) If the reduced form $(a,b,c)$ has the reduced right neighbor $(a',b',c')$, the reduced form $(c,b,a)$ has the left neighbor $(c',b',a')$. If the reduced form $(a,b,c)$ has the left neighbor $({}'a,{}'b,{}'c)$, the reduced form $(c,b,a)$ has the reduced right neighbor $({}'c,{}'b,{}'a)$. The forms $(-{}'a,{}'b,-{}'c)$, $(-a,b,-c)$, $(-a',b',-c')$ are also reduced, with the second a right neighbor of the first, and the third of the second; or the first a left neighbor of the second, and the second of the third. The same holds for the three forms $(-c',b',-a')$, $(-c,b,-a)$, $(-{}'c,{}'b,-{}'a)$. These facts are so clear that no further discussion is needed.
